% ===========================================================================
% ΔΟΚΙΜΑΣΤΙΚΟ ΚΕΦΑΛΑΙΟ ΜΑΘΗΜΑΤΙΚΩΝ
% ===========================================================================
% Η ίδια μακέτα εξυπηρετεί Μαθηματικά και άλλα επιστημονικά βιβλία.
% Δεν απαιτείται διαφορετικό universalbook.sty: το περιεχόμενο αλλάζει.
% Δες παρακάτω τις εντολές definitionbox, theorembox, proof, TikZ και figure.
% ===========================================================================
\BookChapterPhoto{Συναρτήσεις και γραφήματα}{Μαθηματικά — δοκιμή της ίδιας εκδοτικής μακέτας}{assets/math-chapter-art.pdf}{Διανυσματική εικονογράφηση, ενσωματωμένη στο XeLaTeX}{%
\item Να αναγνωρίζεις συναρτήσεις πρώτου βαθμού.
\item Να ερμηνεύεις τους συντελεστές τους.
\item Να συνδέεις τύπους με γραφικές παραστάσεις.
\item Να γράφεις σαφή μαθηματική αιτιολόγηση.
}
\BookSection{Ορισμός και βασική ιδέα}
\section{Η συνάρτηση πρώτου βαθμού}
Ο τύπος \(f(x)=\alpha x+\beta\), με \(\alpha\neq0\), περιγράφει μια συνάρτηση πρώτου βαθμού. Η γραφική της παράσταση είναι ευθεία.
% Ορισμός: αριθμείται αυτόματα ανά κεφάλαιο.
\begin{definitionbox}[Συνάρτηση πρώτου βαθμού]
\(f(x)=\alpha x+\beta\), όπου \(\alpha\neq0\) και \(\alpha,\beta\in\mathbb R\).
\end{definitionbox}

\begin{table}[htbp]
\centering
% v0.23: ο λιτός πίνακας έχει την ΙΔΙΑ σκούρα κεφαλίδα
% και τις ΙΔΙΕΣ αποστάσεις με τον χρωματιστό.
% Μόνη διαφορά: οι γραμμές δεδομένων δεν αποκτούν φόντο.
\begin{BookPlainTable}
\begin{tabularx}{.82\linewidth}{@{}>{\centering\arraybackslash}X>{\centering\arraybackslash}X@{}}
\BookTableHeaderRow
\BookTableHead{$x$} & \BookTableHead{$f(x)=2x+1$}\\
$-1$ & $-1$\\
$0$ & $1$\\
$1$ & $3$\\
\end{tabularx}
\end{BookPlainTable}
\caption{Λιτός πίνακας: ίδιος τίτλος κεφαλίδας, λευκές γραμμές δεδομένων.}
\end{table}

\BookMinorHeading{Παρατήρηση}
Η παράμετρος \(\alpha\) καθορίζει την κλίση της ευθείας, ενώ η παράμετρος \(\beta\) καθορίζει την τεταγμένη του σημείου τομής με τον κατακόρυφο άξονα.
\begin{AttentionBox}
Στο γενικό πρότυπο τα χρώματα σημαίνουν τη λειτουργία ενός πλαισίου, όχι το γνωστικό αντικείμενο. Το ίδιο πλαίσιο «Προσοχή» χρησιμοποιείται με συνέπεια σε Φυσική, Μαθηματικά και άλλα βιβλία.
\end{AttentionBox}
\BookSection{Απόδειξη — υπόδειγμα}
% Θεώρημα: αριθμείται αυτόματα ανά κεφάλαιο.
\begin{theorembox}[Πρόσθεση ίσων ποσοτήτων]
Αν \(a=b\), τότε \(a+c=b+c\) για κάθε πραγματικό αριθμό \(c\).
\end{theorembox}
% Απόδειξη: η ονομασία «Απόδειξη» έρχεται από τον κεντρικό config.
\begin{proof}
Προσθέτοντας τον ίδιο αριθμό στα δύο μέλη της ισότητας, διατηρούμε την ισότητα.
\end{proof}
\begin{workedexample}{Εφαρμογή σε συνάρτηση}
\begin{ProblemBox}
Για τη συνάρτηση \(f(x)=2x+1\), να υπολογιστούν \(f(0)\) και \(f(3)\).
\end{ProblemBox}
\begin{ThinkBox}
Αντικαθιστώ στη μεταβλητή \(x\) την τιμή που μου ζητείται, χωρίς να αλλάζω τον τύπο της συνάρτησης.
\end{ThinkBox}
\SolutionHeading
\[f(0)=2\cdot0+1=1,\qquad f(3)=2\cdot3+1=7.\]
\end{workedexample}
\BookSection{Άσκηση}
\begin{exercise}[MATH-A-001]{}
Να βρεθεί η τιμή του \(x\) για την οποία \(f(x)=9\), όταν \(f(x)=2x+1\).
\Answer{$x=4$.}
\end{exercise}
